% ============================================================
\chapter{TDA and Deep Learning}
\label{ch:deep_tda}
% ============================================================

\begin{intuition}
Integrating TDA into deep learning poses a fundamental challenge:
the computation of persistence is a priori \textbf{not differentiable}.
Recent work has shown how to make persistence differentiable, enabling
end-to-end training of neural networks with topological objectives.
\end{intuition}

% ============================================================
\section{Differentiable persistence}
\label{sec:diff_persistence}
% ============================================================

\begin{definition}[Persistence as a function of filtrations]
\index{differentiable persistence}
Consider a simplicial complex $K$ with a parametric filtration
$f_\theta : K \to \R$ depending on parameters $\theta \in \R^p$.
The persistence diagram $\mathrm{Dgm}(f_\theta)$ is a function
of $\theta$.
\end{definition}

\begin{theorem}[Differentiability --- Gameiro et al.\ 2016, Poulenard et al.\ 2018]
\label{thm:diff_persistence}
\index{differentiability of persistence}
Let $K$ be a finite simplicial complex and $f : K \to \R$ a generic
filtration (distinct filtration values). The birth coordinates
$b_i(\theta)$ and death coordinates $d_i(\theta)$ in the persistence
diagram are differentiable functions of $\theta$ almost everywhere.
More precisely:
\[
  \frac{\partial b_i}{\partial \theta_j} =
  \frac{\partial f_\theta(\sigma_{b_i})}{\partial \theta_j}, \qquad
  \frac{\partial d_i}{\partial \theta_j} =
  \frac{\partial f_\theta(\sigma_{d_i})}{\partial \theta_j}
\]
where $\sigma_{b_i}$ and $\sigma_{d_i}$ are the simplices that
create and kill class $i$.
\end{theorem}

\begin{attention}[title=\textbf{Non-differentiability points}]
Differentiability fails when two filtration values coincide
(the birth and death simplices may change). In practice, one adds
a small perturbation or uses subgradients.
\end{attention}

\begin{keyformulas}[title=\textbf{Persistence gradient}]
\[
  \frac{\partial \mathrm{pers}_i}{\partial \theta}
  = \frac{\partial d_i}{\partial \theta}
  - \frac{\partial b_i}{\partial \theta}
  = \nabla_\theta f(\sigma_{d_i}) - \nabla_\theta f(\sigma_{b_i})
\]
\end{keyformulas}

\begin{minted}{python}
import torch
from topologylayer.nn import RipsLayer

# Differentiable point cloud
X = torch.randn(50, 2, requires_grad=True)

# Differentiable persistence layer
rips_layer = RipsLayer(maxdim=1, sublevel=False)
dgm = rips_layer(X)

# Loss based on H1 persistence
loss = -dgm[1][:, 1].sum()  # maximize H1 persistences
loss.backward()
print(f"Gradient shape: {X.grad.shape}")
\end{minted}

% ============================================================
\section{PersLay: a generic persistence layer}
\label{sec:perslay}
% ============================================================

\begin{definition}[PersLay --- Carriere et al.\ 2020]
\index{PersLay}
\textbf{PersLay} is a neural layer that transforms a persistence
diagram $D = \{(b_i, d_i)\}_{i=1}^n$ into a fixed-dimension vector
via:
\[
  \mathrm{PersLay}(D) = \mathrm{op}\!\left(
  w(p_i) \cdot \phi(p_i)\right)_{i=1}^n
\]
where:
\begin{itemize}
  \item $w : \R^2 \to \R_+$ is a weight function (attenuating
        points near the diagonal).
  \item $\phi : \R^2 \to \R^q$ is a pointwise representation
        function (learned or fixed).
  \item $\mathrm{op}$ is a permutation-invariant aggregation
        (sum, max, top-$k$, attention).
\end{itemize}
\end{definition}

\begin{proposition}[Universality of PersLay]
With suitable choices of $w$, $\phi$, and $\mathrm{op}$, PersLay
can uniformly approximate any continuous, permutation-invariant
function on persistence diagrams.
\end{proposition}

\begin{remark}
PersLay unifies classical representations:
\begin{itemize}
  \item \textbf{Landscapes}: $\phi$ is the tent function,
        $\mathrm{op}$ is top-$k$.
  \item \textbf{Images}: $\phi$ is a Gaussian,
        $\mathrm{op}$ is sum.
  \item \textbf{Silhouettes}: $\phi$ is the tent function,
        $\mathrm{op}$ is the weighted average.
\end{itemize}
\end{remark}

% ============================================================
\section{Topological regularization}
\label{sec:regularization}
% ============================================================

\begin{definition}[Topological loss]
\index{topological regularization}
A \textbf{topological loss} is a term added to the cost function
to encourage or penalize certain topological structures:
\[
  \mathcal{L}_{\mathrm{total}} = \mathcal{L}_{\mathrm{task}}
  + \lambda \cdot \mathcal{L}_{\mathrm{topo}}
\]
where $\mathcal{L}_{\mathrm{topo}}$ is computed from the persistence
diagram.
\end{definition}

\begin{example}[Topological loss for segmentation]
In image segmentation, one can penalize spurious small holes in the
prediction:
\[
  \mathcal{L}_{\mathrm{topo}} = \sum_{(b_i, d_i) \in H_1}
  (d_i - b_i)^2
\]
which forces the network to produce simply connected regions.
\end{example}

\begin{theorem}[Convergence of topological regularization]
Under regularity conditions (generic filtration, Lipschitz),
stochastic gradient descent with topological loss converges to a
critical point of $\mathcal{L}_{\mathrm{total}}$.
\end{theorem}

\begin{minted}{python}
import torch
import torch.nn as nn

class TopologicalLoss(nn.Module):
    """Penalize spurious holes in segmentation."""
    def __init__(self, lam=1.0):
        super().__init__()
        self.lam = lam

    def forward(self, pred, target, dgm_h1):
        ce_loss = nn.functional.cross_entropy(pred, target)
        # H1 persistences = d - b
        pers = dgm_h1[:, 1] - dgm_h1[:, 0]
        topo_loss = (pers ** 2).sum()
        return ce_loss + self.lam * topo_loss
\end{minted}

% ============================================================
\section{Topological autoencoders}
\label{sec:topo_autoencoders}
% ============================================================

\begin{definition}[Topological autoencoder --- Moor et al.\ 2020]
\index{topological autoencoder}
A \textbf{topological autoencoder} is an autoencoder
$(\mathrm{Enc}, \mathrm{Dec})$ whose loss includes a topological
term:
\[
  \mathcal{L} = \norm{x - \mathrm{Dec}(\mathrm{Enc}(x))}^2
  + \lambda \cdot d_W\!\left(\mathrm{Dgm}(X),\;
  \mathrm{Dgm}(\mathrm{Enc}(X))\right)^2
\]
where $d_W$ is a Wasserstein distance between the persistence
diagrams of the input space $X$ and the latent space
$\mathrm{Enc}(X)$.
\end{definition}

\begin{intuition}
The idea is to preserve the \textbf{topology} of the data during
dimensionality reduction. A standard autoencoder (or VAE) may
collapse loops or merge connected components. The topological term
penalizes such distortions.
\end{intuition}

\begin{proposition}[Topological preservation]
If $\mathcal{L}_{\mathrm{topo}} = 0$, then the Betti numbers of
$X$ and $\mathrm{Enc}(X)$ coincide for the Rips filtration used.
\end{proposition}

% ============================================================
\section{Graph neural networks and persistent homology}
\label{sec:gnn_persistence}
% ============================================================

\begin{definition}[Learned graph filtration]
\index{learned filtration}
Let $G = (V, E)$ be a graph with node features $h_v \in \R^d$.
Define a learned filtration:
\[
  f_\theta(v) = g_\theta(h_v), \qquad
  f_\theta(e_{uv}) = \max(f_\theta(u), f_\theta(v))
\]
where $g_\theta : \R^d \to \R$ is a neural network. The persistence
diagram of this filtration provides topological features of the graph.
\end{definition}

\begin{example}[Graph classification]
For molecular graph classification:
\begin{enumerate}
  \item A GNN produces node embeddings $h_v$.
  \item A learned filtration $f_\theta$ is computed.
  \item The persistence diagram is vectorized via PersLay.
  \item The resulting vector is concatenated with the GNN readout.
\end{enumerate}
This combination improves expressiveness beyond the Weisfeiler-Leman
test.
\end{example}

\begin{theorem}[Topological expressiveness --- Horn et al.\ 2022]
\index{expressiveness}
There exist pairs of non-isomorphic graphs indistinguishable by
the $k$-WL test but distinguishable by the persistence of the
filtration induced by a message-passing GNN.
\end{theorem}

% ============================================================
\section{Implementation with PyTorch}
\label{sec:implementation}
% ============================================================

\begin{minted}{python}
import torch
import torch.nn as nn
from torch_geometric.nn import GCNConv, global_mean_pool

class TopoGNN(nn.Module):
    """GNN with topological features."""
    def __init__(self, in_dim, hidden_dim, out_dim):
        super().__init__()
        self.conv1 = GCNConv(in_dim, hidden_dim)
        self.conv2 = GCNConv(hidden_dim, hidden_dim)
        self.filtration = nn.Linear(hidden_dim, 1)
        self.classifier = nn.Linear(hidden_dim + 16, out_dim)

    def forward(self, x, edge_index, batch):
        h = torch.relu(self.conv1(x, edge_index))
        h = torch.relu(self.conv2(h, edge_index))
        # Standard readout
        graph_emb = global_mean_pool(h, batch)
        # Learned filtration
        filt_vals = self.filtration(h).squeeze()
        # ... persistence computation and PersLay ...
        return graph_emb  # simplified
\end{minted}

% ============================================================
\section{Exercises}
% ============================================================

\begin{exercise}
Compute by hand the gradient of the total persistence
$\sum_i (d_i - b_i)$ with respect to filtration values for the
complex $\{a, b, c, ab, bc\}$ with $f(a) = 0, f(b) = 1, f(c) = 2,
f(ab) = 1, f(bc) = 2$.
\end{exercise}

\begin{exercise}
Show that PersLay with $\phi(b, d) = (b, d-b)$, $w \equiv 1$, and
$\mathrm{op} = \mathrm{sum}$ reduces to computing
$(\sum b_i, \sum(d_i - b_i))$.
\end{exercise}

\begin{exercise}[Implementation]
Implement a simple topological autoencoder on a circle-shaped dataset.
Compare the latent space with and without the topological term.
\end{exercise}

\begin{exercise}
Give an example of two non-isomorphic graphs indistinguishable by
$1$-WL but distinguishable by the $H_0$ persistence of a
degree-induced filtration.
\end{exercise}

\begin{exercise}[Project]
Train an image segmentation network (U-Net) with and without
topological regularization on a blood vessel dataset. Compare the
Betti numbers of the predictions.
\end{exercise}
